% !TEX TS-program = lualatex
Протестировать интеллектуальных агентов в игре <<Нарды>> на основе этюдов других игроков достаточно затруднительно, поскольку сила позиции во многом определяется следующим броском игральных костей, который невозможно предсказать заранее. Это делает стандартные методы позиционного анализа менее применимыми по сравнению с играми с полной информацией (отсутствие фактора случайности), такими как шахматы или шашки.

Для оценки динамики обучения интеллектуальных агентов была использована уникальная особенность игры <<Нарды>> --- наличие фиксированного направления движения шашек. Поскольку игрок может перемещать свои шашки только от дома противника в сторону своего собственного, игра имеет направленное течение. Это исключает возможность возникновения циклических ситуаций (при условии равновероятного выпадения всех граней костей), что гарантирует конечность любой партии.

Таким образом, каждый ход, независимо от его эффективности, вносит вклад в продвижение игры к завершению. Это свойство позволяет использовать в процессе обучения простого случайного агента в качестве противника базового уровня сложности. Несмотря на его примитивную стратегию, он обеспечивает достаточно разнообразную игровую среду, в которой можно отслеживать прогресс обучаемого агента и анализировать развитие его игровой стратегии на разных этапах обучения.

Для проведения тестирования были сгенерированы 20 моделей. В \refdf{ai-for-test} представлены их основные характеристики, включая название, количество нейронов в скрытом слое и число сыгранных партий.

\begin{df}{|C|C|C|C|}{m}{Список тестируемых моделей}{ai-for-test}\hline
    Название модели & Число нейронов в скрытом слое & Сыгранное число партий \\ \hline
    TD40v1 & $40$ & $10000$ \\ \hline
    TD40v2 & $40$ & $20000$ \\ \hline
    TD40v3 & $40$ & $30000$ \\ \hline
    TD40v4 & $40$ & $40000$ \\ \hline
    TD40v5 & $40$ & $50000$ \\ \hline
    TD40v6 & $40$ & $60000$ \\ \hline
    TD40v7 & $40$ & $70000$ \\ \hline
    TD40v8 & $40$ & $80000$ \\ \hline
    TD40v9 & $40$ & $90000$ \\ \hline
    TD40v10 & $40$ & $100000$ \\ \hline
    TD80v1 & $80$ & $10000$ \\ \hline
    TD80v2 & $80$ & $20000$ \\ \hline
    TD80v3 & $80$ & $30000$ \\ \hline
    TD80v4 & $80$ & $40000$ \\ \hline
    TD80v5 & $80$ & $50000$ \\ \hline
    TD80v6 & $80$ & $60000$ \\ \hline
    TD80v7 & $80$ & $70000$ \\ \hline
    TD80v8 & $80$ & $80000$ \\ \hline
    TD80v9 & $80$ & $90000$ \\ \hline
    TD80v10 & $80$ & $100000$ \\ \hline
\end{df}

В качестве метрик оценки качества моделей были выбраны следующие показатели:

\begin{enumerate}
    \item Суммарное число выигранных матчей против моделей предыдущих поколений. Эта метрика отражает прогресс модели относительно её предшественников и позволяет оценить динамику развития в ходе обучения.
    \item Число побед над случайным агентом. Данный показатель демонстрирует, насколько сложной и эффективной является модель по сравнению с примитивным, не обученным противником, что особенно важно при оценке её применимости к реальным задачам.
\end{enumerate}

Для оценки производительности моделей в каждом эксперименте было проведено по 900 партий между различными поколениями одной и той же модели, а также 1000 партий с участием случайного агента.

На \refimage{exp1} представлены результаты для 10 поколений модели версии T40, в архитектуре которой скрытый слой содержит 40 нейронов. График напоминает параболу с направленными вбок ветвями. Можно наблюдать устойчивую положительную динамику: с каждым новым поколением модель приобретала новые знания и улучшала результаты по сравнению с предыдущими версиями. Это говорит о стабильном процессе обучения и адаптации к игровой среде.

\image[width=0.8\textwidth]{Победы интеллектуального агента TD40}{exp1}{exp1}

На \refimage{exp3} визуализированы интеллектуальные способности модели в сравнении со случайным агентом. Уже после 20000 сыгранных партий модель демонстрирует уровень игры, превышающий базовый, и стабильно выигрывает у случайного агента. Наилучшее значение зафиксировано у последнего, десятого поколения: оно выигрывает примерно 6 партий из 10, что указывает на наличие устойчивой стратегии.

\image[width=0.8\textwidth]{Победы интеллектуального агента TD40 против случайного агента}{exp3}{exp3}

На \refimage{exp2} показана динамика развития моделей типа TD80, в которых используется скрытый слой из 80 нейронов. В отличие от графика модели T40, данный график имеет форму изогнутой кривой с ярко выраженным изломом на шестом поколении. Такой излом может свидетельствовать о начале переобучения --- модель начинает запоминать специфические сценарии вместо того, чтобы обобщать информацию, что снижает её способность к адаптации.

\image[width=0.8\textwidth]{Победы интеллектуального агента TD80}{exp2}{exp2}

На \refimage{exp4} представлена сложность поведения моделей по результатам матчей со случайным агентом. Обратим внимание, что модель второго поколения не занимает лидирующую позицию, несмотря на раннее развитие. В то же время модели, начиная с четвёртого поколения, стабильно демонстрируют уверенную игру и в среднем выигрывают у случайного агента 7 партий из 10, что говорит о возросшей сложности и эффективности поведения.

\image[width=0.8\textwidth]{Победы интеллектуального агента TD80 против случайного агента}{exp4}{exp4}

Таким образом, на основе полученных данных можно сделать следующие выводы:
\begin{enumerate}
    \item Лучшими результатами в обученной выборке обладает модель TD80v3. Несмотря на незначительный излом в развитии по сравнению со вторым поколением, она демонстрирует уверенное превосходство над случайным агентом и стабильно выигрывает большинство матчей.
    \item Модели с увеличенным числом нейронов на скрытом слое (80 нейронов) показывают более быстрое обучение, но также подвержены раннему переобучению, которое, как правило, начинается с шестого поколения.
    \item В отличие от них, модели типа T40, имеющие меньшее число нейронов (40), демонстрируют более стабильное и устойчивое обучение. В рамках данного эксперимента признаки переобучения для них не были зафиксированы даже на поздних этапах.
\end{enumerate}
